\chapter{Controles de vuelo}
\label{cap:mm-controles}

\opening{Tubos de empuje, palancas acodadas y un plato oscilante: cada mando del R22 llega a los rotores por un camino mecánico que hay que armar, frenar y verificar sin errores.}

\section{Descripción}
\fuente{MM, capítulo 8}

Los mandos dobles son equipo estándar. Todos los controles primarios se accionan mediante \textbf{tubos de empuje-tracción (\emph{push-pull})} y \textbf{palancas acodadas (\emph{bellcranks})}. Los rodamientos del sistema son de bolas sellados o de camisa autolubricada.

El cíclico está montado al centro: las empuñaduras izquierda y derecha van en un tubo transversal que pivota sobre el bastón central. En aeronaves recientes el mecánico puede ajustar el ángulo de la empuñadura del piloto hacia adelante o atrás respecto del tubo transversal. El colectivo es convencional, con acelerador de empuñadura giratoria; al subir el colectivo, una articulación de interconexión abre el acelerador y un gobernador electrónico hace las correcciones finas para mantener las RPM. En aeronaves recientes, una empuñadura fija de plástico detrás de la giratoria permite al piloto apoyar la mano sin interferir con el gobernador.

\begin{advertencia}
El armado de los controles de vuelo es crítico y requiere inspección por una persona calificada. Si no hay una segunda persona disponible, RHC recomienda que el instalador tome un \textbf{descanso de 5 minutos} antes de inspeccionar las conexiones que él mismo armó.
\end{advertencia}

En todas las uniones con rótula de los controles se usan separadores A115-1 a cada lado de la bola y, donde se indica, arandelas de seguridad A214-3 bajo la cabeza del perno: si la rótula se rompe, la arandela de mayor diámetro impide que el extremo del tubo se desprenda. Salvo indicación, los pernos se aprietan con torque estándar (\S\,23-32), se instala contratuerca de chapa (\emph{palnut}) y se marca con pintura de torque (figura 2-1 del MM); véase el capítulo~\ref{cap:mm-practicas}.

\section{Cíclico}

\subsection{Conjunto del cíclico}

Para retirar el cíclico se retiran primero el colectivo, los amortiguadores C683 si los hay, las tapas de los cables de mezcla y calefacción del carburador, las perillas de fricción y de compensación derecha, y la tapa de la caja del cíclico (desconectando sus conectores). Luego se desconectan tubos y horquilla, se lleva el cíclico a fondo a la derecha para aliviar el resorte de compensación, se aplica fricción para mantener armada la pila de fricción y se levanta la caja recta entre los paneles de quilla.

\begin{precaucion}
Debe verificarse el reglaje (capítulo~\ref{cap:mm-reglaje}) si:
\begin{itemize}
\item se reemplazó el conjunto del cíclico, la caja A338-1 o el bastón A320-1;
\item se reemplazó o modificó la horquilla A205-1 (une colectivo y pivote del cíclico);
\item cambió la distancia entre centros de rótulas de los tubos A121-1, -3, -5 o -7;
\item cambió la longitud de los soportes del eje intermedio (\emph{jackshaft}).
\end{itemize}
\end{precaucion}

Datos de la instalación:
\begin{itemize}
\item \textbf{Horquilla A205-1}: 3,50--3,60~in del centro de la rótula al eje del agujero del perno (sin palnut en esta rótula). Si cambia esta medida, verifique las RPM de autorrotación (\S\,10.250).
\item Con el cíclico adelante, el tubo A121-3 se une a las palancas A958-1 (o A175-1); se admite un mínimo de 1 hilo visible más allá de la tuerca.
\item \textbf{Holgura de los controles}: mínimo \textbf{0,12~in} en todo el recorrido, sin trabas ni interferencias, salvo:
 \begin{itemize}
 \item 0,030~in mínimo entre los tubos A121-7 y los pernos del bastidor superior;
 \item contactos permitidos: extremos inferiores de A121-3/-7 en las esquinas traseras; A121-1 con el resorte de sobrecarrera A327-1 (adelante a la derecha); A121-3 con la biela del colectivo B328-1 con colectivo arriba y cíclico atrás a fondo.
 \end{itemize}
\end{itemize}

\subsection{Palancas acodadas A958-1 / A175-1}

Las palancas A175-1 de los números de serie 0001--0356 (excepto 0256, 0301, 0350 y 0351) tienen aplicabilidad limitada; con el soporte B379 se requiere la palanca A958-1.

\subsection{Empuñadura del cíclico}

El ángulo de la empuñadura se ajusta girándola a la posición deseada y apretando los tornillos de la tapa D379-1 con torque especial de \textbf{40~in-lb}. Verifique holgura del cableado y ausencia de trabas.

\subsection{Resorte de compensación derecha}

Compensa la tendencia del cíclico a irse a la izquierda en vuelo. Resortes disponibles:

\begin{center}
\begin{tabularx}{\linewidth}{@{}llX@{}}
\toprule
Resorte & Constante & Suplementos \\
\midrule
A618-1 & 10,7~lb/in & máx.\ 3 arandelas NAS1149F0463P \\
A618-2 (negro) & 14,5~lb/in & no se suplementa \\
\bottomrule
\end{tabularx}
\end{center}

Los conjuntos A056 (S/N 0002--0549) deben consultarse con RHC.

\subsection{Cordón elástico (compensación adelante)}

Los cordones A918-1 a -13 contrarrestan la fuerza del cíclico hacia atrás. Si en crucero la empuñadura tiende a ir hacia adelante, use un cordón más débil; si tiende a ir atrás, uno más fuerte. El cordón corto se engancha entre la placa A697-2 y la lengüeta A697-6 del extremo delantero del tubo A121-1; el largo se pasa sobre el soporte A315-1 y se unen sus ganchos. Verifique en todo el recorrido la holgura de los ganchos con el tubo, el mazo de cables y el cableado del cíclico.

\subsection{Fricción del cíclico}

\begin{itemize}
\item Perilla: entre 1/8 y 1 vuelta libre antes de empezar a dar fricción.
\item Con fricción desactivada: fuerza media longitudinal \textbf{$\le$ 3~lb}; lateral \textbf{$\le$ 2~lb}.
\item Con fricción total, fuerza lateral \textbf{$\le$ 10~lb}.
\item Pila de fricción: arandelas Belleville grandes de 1,1~in (las antiguas de menor diámetro, 4 unidades, se reemplazan por A478-1); separador A130-18 de 0,105--0,115~in; A130-19 de 4,47~in; cota de montaje 0,035--0,055~in (los separadores se lijan para lograrla).
\end{itemize}

\subsection{Tubo de empuje A121-7}

\textbf{Inspección y reparación}:
\begin{itemize}
\item Rayas longitudinales de hasta 0,010~in de profundidad, o transversales de hasta 0,005~in, pueden pulirse.
\item Abolladuras de más del 5\,\% del diámetro: reemplazar el tubo.
\item Desgaste máximo en la guía: \textbf{0,015~in por pared} después de pulir; más allá, reemplazar. Todo tubo con desgaste requiere manguito.
\end{itemize}

\textbf{Manguito}: el adhesivo se limita con cinta a una zona de 6,5~in situada a 31~in del extremo del tubo (no de la rótula). Pruebe primero una banda de 1~in: si la pintura se ampolla, quítela de toda la zona. El manguito debe instalarse \textbf{dentro de los 2 minutos} de aplicado el adhesivo y se contrae con pistola de calor de \textbf{máx.\ 1200~W}. Rechace el manguito si tiene arrugas, poros, burbujas, cortes o zonas rotas.

\textbf{Instalación}:
\begin{pasos}
\item Inspeccione la guía; reemplácela si está desgastada o hay frotamiento (\emph{fretting}) entre el material y la placa móvil.
\item Las guías antiguas remachadas al fuselaje deben reemplazarse por la guía móvil (kit KI-19A, SB-36) si los tubos A121-7 están desgastados.
\item Monte la guía móvil sin apretar; instale el tubo con extremo cuidado para no dañar el manguito, con los separadores correctos; pernos a 100~in-lb + arrastre de tuerca, también en el plato oscilante.
\item Suba el colectivo unas 7~in (empuñadura del acelerador al panel central) con el cíclico neutro; apriete los tres pernos de la guía hasta leve fricción, verifique que los tubos queden centrados y aplique \textbf{40~in-lb}.
\item Verifique holguras en todo el recorrido: pernos delanteros del bastidor superior (cortafuego horizontal) a los tubos A121-7, mín.\ \textbf{0,010~in}; ala interior superior del tanque de combustible al tubo izquierdo, mín.\ \textbf{0,10~in}.
\item Si se desenroscaron rótulas, verifique el reglaje del rotor principal.
\end{pasos}

\section{Colectivo}

\subsection{Bastón del colectivo}

\begin{advertencia}
El resorte del colectivo está comprimido. Antes de soltarlo, con el colectivo abajo, instale el retén MT294-1 o átelo con alambre de frenado de 0,032~in dado varias vueltas de rótula a rótula. Al terminar, retire el retén o el alambre: olvidarlo puede limitar el recorrido del control y crear un peligro para la seguridad del vuelo.
\end{advertencia}

Puntos clave de la instalación:
\begin{itemize}
\item Separador A115-1 con el diámetro menor contra el bloque de rodamientos A329-1.
\item Tope del colectivo A333-1 al anclaje de cinturón A348-1 con arandelas que dejen un \textbf{juego axial de 0,001--0,035~in} sin trabas.
\item Resorte de sobrecarrera A327-1 al brazo del acelerador: tornillo A486-1 con brazo C342-1, A486-4 con C341-3; cabeza del tornillo y diámetro menor del separador A341-1 contra la rótula. Una instalación incorrecta puede trabar o separar la rótula y hacer perder el control del acelerador.
\item Reglaje de correlación del acelerador (\S\,10.150) y ajuste del resorte según vuelo de prueba.
\end{itemize}

\subsection{Resorte del colectivo}

Equilibra las fuerzas del colectivo en vuelo. Se ajusta o se cambia el resorte y se prueba en vuelo de crucero nivelado hasta obtener las fuerzas deseadas.

\begin{center}
\small
\begin{tabular}{@{}llll@{}}
\toprule
Conjunto & Resorte & Color & Constante \\
\midrule
A038-1 & A429-1 (antiguo) & plateado & 89~lb/in \\
A038-2 & A429-2 (antiguo) & imprimación & 35~lb/in \\
A038-3 o -5 & A429-3 & gris & 50~lb/in (débil) \\
A038-4 & A429-4 (antiguo) & negro & 70~lb/in \\
A038-6 & A429-5 & dorado & 50~lb/in (estándar) \\
A038-7 & A429-8 & blanco & 71~lb/in (el más fuerte) \\
\bottomrule
\end{tabular}
\end{center}

A038-5 reemplaza a A038-2; A038-6 reemplaza a A038-1, -3 y -4. Para instalarlo, el conjunto se comprime a \textbf{4,00--4,30~in} entre centros de rótulas, con el extremo de tuerca de las guías hacia arriba. En los conjuntos A038-5, -6 y -7 \textbf{ambas rótulas deben estar a fondo} (la B292-3 es de rosca izquierda); de lo contrario una puede quedarse sin rosca durante el ajuste.

\textbf{Ajuste}: enroscar la rótula en la tapa \emph{disminuye} la fuerza hacia arriba; extenderla la \emph{aumenta}. En A038-1 a -4 la extensión máxima de la rótula inferior es 1,1~in de la tapa al centro de la rótula. En A038-5/-6/-7 se gira el resorte a mano con el colectivo arriba (rótula inferior rosca derecha, superior izquierda). Verifique siempre que las espiras no se toquen con el colectivo abajo.

Para cambiar el resorte se usa una morsa de mordazas blandas; inspeccione las varillas guía A428-1 (reemplazar si tienen rayas o mellas) y engráselas levemente con A257-1.

\subsection{Fricción del colectivo}

Arandelas A478-2 (o B0500-022-S) con las caras cóncavas enfrentadas; arandelas A141-20 con la cara de teflón gris contra el tope A333-1.

\begin{center}
\begin{tabular}{@{}lc@{}}
\toprule
Condición (en la empuñadura, tirando hacia arriba) & Fuerza \\
\midrule
Fricción bloqueada & 12--20~lb \\
Fricción desbloqueada & 6--10~lb \\
\bottomrule
\end{tabular}
\end{center}

\begin{advertencia}
Una fricción del colectivo mayor de 20~lb puede impedir entrar en autorrotación.
\end{advertencia}

\section{Gobernador de RPM}

El gobernador mantiene las RPM del motor detectando variaciones y aplicando correcciones al acelerador a través de un \textbf{embrague de fricción} que el piloto puede vencer fácilmente. Está activo solo por encima del 80\,\% de RPM del motor y se conecta con el interruptor del extremo del colectivo derecho. Ayuda a controlar las RPM en condiciones normales pero puede no impedir sobrevelocidades o bajas RPM en maniobras agresivas.

\begin{precaucion}
A gran altitud de densidad la respuesta del gobernador puede ser demasiado lenta para evitar una sobrevelocidad con ráfagas, recogidas o al bajar el colectivo.
\end{precaucion}

\textbf{Controlador}: el D270-1 (con unidad de monitoreo del motor, EMU; \S\,33-137) va en el mamparo derecho detrás del respaldo derecho; el B286-2, en el respaldo izquierdo. Antes de trabajar, batería apagada y disyuntor GOV (2~A) afuera. \textbf{No tiene ajustes externos}: si falla, se devuelve a RHC.

\textbf{Conjunto del gobernador} (en el colectivo): no se permite ajustar el embrague de fricción ni reemplazar el motorreductor en campo; si funcionan mal, se devuelve a RHC. Biela de enlace: \textbf{2,47--2,53~in} entre centros; brazo de salida hacia arriba, hacia la biela B328-1. Con tornillos AN503-8-4: 27~in-lb y frenado con alambre de 0,020~in.

\subsection{Diagnóstico}

La mayoría de los problemas del gobernador se deben a los \textbf{contactos del tacómetro del magneto derecho} del motor (lado izquierdo del helicóptero) desajustados o defectuosos (véase el manual de servicio de TCM). Entre 80 y 115\,\% el gobernador trata de mantener \textbf{104\,\% $\pm$ 1,5\,\%} (banda muerta de 3\,\%, 102,5--105,5\,\%). Para medirla:

\begin{pasos}
\item Con una persona habilitada a los mandos, arranque y haga la prueba de motor según el POH, con el colectivo totalmente abajo.
\item Sujetando suavemente el acelerador, suba muy despacio las RPM (sin pasar 107\,\%) y anote dónde aparece la resistencia del gobernador.
\item Baje muy despacio (sin bajar de 97\,\%) y anote el punto de resistencia.
\item La diferencia debe ser de alrededor de 3\,\%.
\end{pasos}

\begin{itemize}
\item 3\,\% pero centrada fuera de 104\,\% $\pm$ 0,5\,\%: problema del controlador.
\item Más de 3\,\% centrada en 104\,\%: exceso de fricción en el varillaje o falta de fricción del gobernador.
\item \textbf{Fricción del varillaje del acelerador}: con el resorte de sobrecarrera desconectado del brazo C342 y una balanza, máx.\ \textbf{4~lb} en movimiento hasta acelerador pleno. Causas de exceso: rótulas trabadas, interferencias, bujes del eje del acelerador ovalizados, bomba de aceleración trabada.
\item \textbf{Fricción del gobernador}: colectivo abajo, fricción aplicada, brazo C342 horizontal, balanza tangente. Fricción en movimiento \textbf{mínimo 8~lb}; la de arranque suele ser 0--0,5~lb mayor; si es 1~lb o más mayor, el embrague puede estar dañado o contaminado.
\item Se requiere una relación mínima \textbf{2:1} entre fricción del gobernador y fricción del varillaje.
\item Funcionamiento errático: contactos del tacómetro (contaminación por exceso de lubricación del fieltro, fuga de aceite por el sello, humedad por el tapón de venteo; oxidación) o daño del cableado blindado.
\item En turbulencia o con el motor poco cargado es normal que fluctúe la presión de admisión; cualquier holgura en el varillaje del acelerador produce oscilaciones de RPM y MAP.
\end{itemize}

\section{Eje intermedio (jackshaft) y montantes}

No hace falta verificar el reglaje si no se aflojaron las rótulas de los montantes ni de los tubos.

\begin{itemize}
\item Montaje del eje intermedio: arandela de seguridad A214-3 entre rótula del montante trasero y cabeza del perno; cabezas de los pernos superiores hacia atrás; 100~in-lb + arrastre, palnut.
\item Tubos A121-3 y -5: cabezas de perno hacia adelante; la rótula delantera del A121-3 lleva A214-3 y separador A115.
\item Montantes: no desarme el montante horizontal largo del vertical; mida y rotule la longitud del montante delantero antes de retirarlo. Dos arandelas AN960-516L en los soportes delanteros. Verifique el enganche de las rótulas por el agujero testigo.
\item Si se modificó un montante, ajuste para que el eje intermedio quede nivelado con el panel de quilla dentro de \textbf{$\pm$0,5°}.
\end{itemize}

\section{Plato oscilante}

Para retirarlo se retiran el rotor principal, las mitades del yugo A203, los topes de caída y de balanceo, la funda; se desconectan los tres tubos y la horquilla A205 inferior y se levanta del tubo deslizante.

\begin{precaucion}
Las bolas plásticas del plato oscilante son obsoletas y deben reemplazarse.
\end{precaucion}

Puntos clave de la instalación:
\begin{itemize}
\item Verifique que el plato gire libremente.
\item Bieleta de paso en la oreja opuesta al interruptor: cinco arandelas AN970-4 y dos AN960-416L bajo la cabeza, A115-1 a cada lado de la rótula, la pila de arandelas por fuera de la oreja; 2 a 4 hilos visibles tras el torque.
\item Oreja del lado del interruptor: bieleta con A214-3 bajo la cabeza; luego horquilla A205 (vástago de rótula hacia arriba).
\item Uniones horquilla/brazo y yugo/brazo: con bujes C648-2 (marrón oscuro), suplementos A117-45/-46/-47 para un \textbf{juego axial de 0,001--0,010~in}; fricción de pivote \textbf{máx.\ 2,5~in-lb}. Con bujes plásticos A139-1, escariado y refrentado (véase abajo).
\item Yugo A203 con chaveta A210-1 (borde redondeado en la ranura del eje); apriete primero el perno de la chaveta. Una pequeña separación entre mitades del lado del brazo de cuerda es normal.
\item Contrapesos A255 y arandelas AN970-4 en el brazo del yugo según balanceo.
\item Plato inferior: horquilla A205 a la izquierda de la oreja central del grupo de tres; tubos delanteros a la izquierda de las otras dos; tubo trasero a la derecha de la oreja trasera. A214-3 entre cabeza y A115-1; A115-1 a cada lado de todas las rótulas.
\item Topes de balanceo orientados según la figura 26-5. Al final: sin trabas en todo el recorrido; instalar rotor, hacer \emph{tracking} y balanceo (\S\,10.200).
\end{itemize}

\textbf{Plato con bola de aluminio y manguito superior sin brida}: se mide el espacio entre el manguito esférico superior y el plato inferior y se elige un separador A209 (incrementos de 0,005~in) que lo llene. Escudo con tornillos a 17~in-lb. La fuerza para inclinar el plato, medida con balanza en el agujero del brazo, debe ser \textbf{libre sin holgura hasta 5~lb} en movimiento (no use la lectura de arranque). Un separador más grueso aumenta la fuerza.

\section{Controles del rotor de cola}

\subsection{Pedales}

Los bloques de rodamiento superior e inferior de los pedales son \textbf{pares mecanizados juntos} (letra o número coincidente): no los mezcle. El pedal derecho va adelante en los bloques y el izquierdo atrás. Ranuras de los bloques llenas de grasa A257-1; pernos de los bloques a 40~in-lb + arrastre. Si hacen falta más de \textbf{5~lb} para mover los pedales, verifique el emparejamiento de los bloques. Tubos A121-9 (palanca A317-1 superior--pedal izquierdo) y A121-11 (inferior--pedal derecho), 100~in-lb + arrastre; las cabezas de los pernos en los pedales, enfrentadas.

\subsection{Palancas acodadas}

\begin{itemize}
\item \textbf{A317-1 inferior delantera}: agujero de pasador de reglaje de 3/16~in arriba del pivote, palanca detrás del soporte. Fuerza máx.\ \textbf{2~lb} con balanza para moverla; si es mayor, escariar bujes (\S\,8.600). Reglaje de tubos: pedales juntos en neutro con un separador de \textbf{2,90~in} entre cada pedal interior y su tope; luego, pedal izquierdo a fondo y ajuste el A121-13 para \textbf{0,06~in mínimo} entre la palanca A013-3 y el cortafuego vertical (se permite tocar la espuma acústica). Perno del A121-13 con la cabeza a la derecha; arandela AN960-416 entre la rótula y la placa izquierda.
\item \textbf{A013-3 superior}: escariador de 0,375/0,376~in si hace falta; arandelas A141-3 junto a los bujes A139-1 azules. Cabeza del perno del A121-13 hacia la derecha.
\item \textbf{A331-1 intermedia}: arandela A141-3 a cada lado; si falta holgura del perno del A121-17 con la voluta, dos AN960-416 bajo la palanca con perno NAS1304-40; A214-3 bajo las cabezas.
\item \textbf{A120-1 trasera y control de paso A031}: arandelas A117-20/-21/-22 entre palanca y espárrago A125-1 para llenar el espacio sin precargar el rodamiento (\textbf{$\pm$0,003~in}); tuerca del espárrago a 70~in-lb + arrastre sujetando el espárrago; pivote a 100~in-lb + arrastre.
\end{itemize}

\subsection{Bieletas de paso del rotor de cola}

\begin{advertencia}
Ambas bieletas deben tener el mismo número de parte. Está prohibido mezclar bieletas de una pieza con ajustables, o de una pieza de acero con de aluminio. Lleve arandelas A214-3 junto a la rótula.
\end{advertencia}

Marque cada bieleta con su pala y oreja antes de retirarla. Una bieleta ajustable nueva se arma a la longitud nominal de \textbf{2,360~in} entre centros; la rótula hembra va por fuera del brazo de control de paso y la macho al cuerno de la pala (se pueden poner arandelas bajo la tuerca para balanceo de cuerda). Luego verifique el reglaje del rotor de cola y haga el balanceo dinámico (capítulo~\ref{cap:mm-reglaje}).

\section{Bujes del yugo A203 y la horquilla A205}

Hay tres tipos de bujes; puede haber más de uno en una instalación, pero \textbf{cada eje de pivote debe tener el mismo tipo}:

\begin{center}
\small\begin{tabularx}{\linewidth}{@{}>{\raggedright\arraybackslash}p{0.2\linewidth}>{\raggedright\arraybackslash}p{0.25\linewidth}>{\raggedright\arraybackslash}X@{}}
\toprule
Buje & Alojamiento & Instalación \\
\midrule
C648-2 (camisa de teflón) & 0,5625/0,5630~in & Imprimación de cromato de zinc o epoxi húmeda; prensar. Suplementos para juego axial. \\
A139-1 (Turcite azul) & 0,499/0,501~in & Prensar con imprimación húmeda; escariar a 0,375~in (o 0,376--0,377) hasta leve arrastre con muñones A105-11; refrentar con herramienta de 1~in y piloto de 0,25~in hasta que el ancho sea 0,005~in menor que la distancia entre arandelas A141-3. Perno a 100~in-lb + arrastre. \\
B115-1 esférico, interferencia & 0,6239/0,6244~in & Calentar a $\approx$180~°F; uno con cromato de zinc a fondo, el otro con grafito (Lubri-Kote) a 0,010--0,020~in sobre el borde. \\
B115-1 esférico, deslizante & 0,6255/0,6260~in (máx.) & Loctite 271 en la pista exterior, sin que entre al rodamiento; armar con sellador húmedo y esperar el curado. \\
\bottomrule
\end{tabularx}
\end{center}

\textbf{Fuerza de pivote del conjunto armado: máx.\ 5~in-lb} (llave de torque en la cabeza del perno). Para reducirla: con A139-1, escariar levemente sobremedida; con B115-1 de interferencia, golpear suavemente la pista exterior del rodamiento instalado con grafito con un tubo de 7/16~in para quitar la precarga axial (nunca girarlo con taladro: daña el teflón); si no basta, cambiar a un conjunto con rodamiento deslizante.
